\documentclass[11pt]{article}
\usepackage[margin=1in]{geometry}
\usepackage{amsmath,amssymb,amsthm,hyperref}
\hypersetup{hidelinks}
\newtheorem{theorem}{Theorem}
\newtheorem{proposition}[theorem]{Proposition}
\theoremstyle{remark}\newtheorem{remark}[theorem]{Remark}
\newcommand{\HS}{\operatorname{HS}}
\newcommand{\pd}{\operatorname{pd}}
\newcommand{\G}{\mathcal G}
\title{Componentwise polymatroidality is not preserved by homological shifts}
\author{}
\date{September 8, 2026}
\begin{document}
\maketitle
\begin{abstract}
We give a negative answer to a question of Ficarra on homological shifts of componentwise polymatroidal ideals. Over any field, the ideal $(x_1,\ldots,x_r,yz,y^a)$ is componentwise polymatroidal for every $r\geq1$ and $a\geq3$, but each of its intermediate homological shift ideals fails polymatroidality in every sufficiently large homogeneous component. A tensor product of two elementary resolutions gives all the shifts, and a single exchange obstruction proves the failure uniformly. The smallest example also contradicts the first-shift assertion in Theorem 10.1 of Ficarra--Lu, arXiv:2509.11977v1; we identify the invalid step in its proof.
\end{abstract}

\section{Statement and context}
Let $R$ be a polynomial ring over a field $K$, with its standard grading. For a monomial ideal $J$, write $J_{\langle d\rangle}$ for the ideal generated by its degree-$d$ monomials. The ideal $J$ is \emph{componentwise polymatroidal} if each nonzero $J_{\langle d\rangle}$ is polymatroidal. Recall that a monomial ideal generated in one degree is polymatroidal if, whenever $u,v$ are minimal generators and $\deg_tu>\deg_tv$, there is a variable $s$ with $\deg_su<\deg_sv$ such that $su/t$ is a generator.

The $i$th homological shift ideal $\HS_i(J)$ is generated by the monomials corresponding to the multigraded shifts in homological position $i$ of a minimal multigraded free resolution of $J$. Thus $\HS_0(J)=J$; the indexing refers to the ideal, rather than its quotient.

Ficarra \cite[Question 9]{F} asks whether every homological shift of a componentwise polymatroidal ideal is again componentwise polymatroidal. For equigenerated polymatroidal ideals, preservation of polymatroidality has been proved by Cid-Ruiz, Matherne, and Shapiro \cite[Theorem A]{CMS}. Ficarra--Lu \cite[Theorem 10.1]{FL} assert the first-shift case of the componentwise question. The following family answers the question negatively and shows that this latter assertion, as stated in the cited version, is false.

\begin{theorem}\label{main}
Let $K$ be any field, $r\geq1$, and $a\geq3$. In
\[
 R=K[x_1,\ldots,x_r,y,z],\qquad I=(x_1,\ldots,x_r,yz,y^a),
\]
the ideal $I$ is componentwise polymatroidal and $\pd_R I=r+1$. For every $1\leq i\leq r$ and every $d\geq a+i$, the ideal $\HS_i(I)_{\langle d\rangle}$ is not polymatroidal. The top shift is the principal ideal
\[
 \HS_{r+1}(I)=(x_1\cdots x_r y^a z).
\]
\end{theorem}

In particular, failure occurs in arbitrarily many homological positions simultaneously, and it persists in all sufficiently large components of each intermediate shift. The proof does not depend on the characteristic of $K$.

\section{Proof}
Put $X=\{x_1,\ldots,x_r\}$. For an integer $j$, let $E_j(X)$ be the ideal generated by the squarefree monomials of degree $j$ in $X$, with $E_0(X)=R$ and $E_j(X)=0$ for $j<0$ or $j>r$.

First we verify the hypothesis of Theorem~\ref{main}. The degree-one component of $I$ is $(X)$. For $2\leq d<a$, a degree-$d$ monomial is in $I$ precisely when it involves a variable of $X$ or both $y$ and $z$. Therefore
\[
 \G(I_{\langle d\rangle})=
 \{u:\deg u=d,\ \deg_yu\leq d-1,\ \deg_zu\leq d-1\}.
\]
For $d\geq a$, the only degree-$d$ monomial outside $I$ is $z^d$, so
\[
 \G(I_{\langle d\rangle})=\{u:\deg u=d,\ \deg_zu\leq d-1\}.
\]
Both descriptions are bounded Veronese sets and satisfy exchange directly: if $\deg_tu>\deg_tv$, equality of total degrees supplies a variable $s$ with $\deg_su<\deg_sv$, and increasing $\deg_su$ by one respects every upper bound since $\deg_su+1\leq\deg_sv$. Thus $I$ is componentwise polymatroidal.

\begin{proposition}\label{shifts}
For every $i\geq0$,
\begin{equation}\label{formula}
 \HS_i(I)=E_{i+1}(X)+yzE_i(X)+y^aE_i(X)+y^azE_{i-1}(X).
\end{equation}
\end{proposition}
\begin{proof}
Write $A=K[X]$ and $B=K[y,z]$. There is an isomorphism
\[
 R/I\cong A/(X)\otimes_K B/(yz,y^a).
\]
The first factor has its minimal Koszul resolution, with squarefree $x$-multidegrees of degree $j$ in position $j$. The second has the minimal multigraded resolution
\[
 0\longrightarrow B(-y^az)
 \xrightarrow{\binom{y^{a-1}}{-z}}
 B(-yz)\oplus B(-y^a)
 \xrightarrow{(yz\ \ y^a)} B
 \longrightarrow B/(yz,y^a)\longrightarrow0.
\]
The notation $B(-u)$ records the multidegree of the monomial $u$. Exactness follows because a relation $fyz+gy^a=0$ gives $fz=-gy^{a-1}$, and the coprime monomials $z,y^{a-1}$ divide $g,f$, respectively. All matrix entries have positive degree. Tensoring these resolutions over the field $K$ gives a minimal free resolution over $R$: exactness follows from the tensor product over a field, and minimality follows because all differential entries remain in the homogeneous maximal ideal. Multiplying the multidegrees and adding homological positions gives the shifts of $R/I$. Position $i+1$ of that resolution is position $i$ of the resolution of $I$, yielding \eqref{formula}.
\end{proof}

For $i=r+1$, only the last summand of \eqref{formula} survives and is nonzero; for $i>r+1$ all summands vanish. This proves the projective-dimension and top-shift assertions.

Fix $1\leq i\leq r$, choose $F\subseteq X$ with $|F|=i-1$, and choose $x_t\in X\setminus F$. Write $x_F=\prod_{x\in F}x$. For $q\geq0$, set
\[
 u=x_Fx_t y^{a-2}z^{q+2},\qquad
 v=x_Fy^az^{q+1}.
\]
These monomials have degree $d=a+i+q$ and belong to $\HS_i(I)$: the second summand of \eqref{formula} contains $u$, and the fourth contains $v$. They are consequently minimal generators of $\HS_i(I)_{\langle d\rangle}$. The exponent of $x_t$ in $u$ exceeds that in $v$, and the only variable with smaller exponent in $u$ than in $v$ is $y$. Exchange would therefore require
\[
 w=yu/x_t=x_Fy^{a-1}z^{q+2}
\]
to belong to $\HS_i(I)$. It does not: its $x$-support has size $i-1$, excluding the first three summands of \eqref{formula}, while its $y$-exponent is smaller than $a$, excluding the fourth. This proves Theorem~\ref{main}.

\section{The smallest example and the first-shift assertion}
For $r=1$ and $a=3$, the example is
\[
 I=(x,yz,y^3),\qquad \HS_1(I)=(xyz,xy^3,y^3z).
\]
Its degree-four component contains $xyz^2$ and $y^3z$ but not the exchange monomial $y^2z^2$. Meanwhile $I_{\langle2\rangle}=(x^2,xy,xz,yz)$, and for every $d\geq3$ the degree-$d$ monomials of $I$ are all such monomials except $z^d$.

The proof of \cite[Theorem 10.1]{FL} asserts
\begin{equation}\label{invalid}
 \HS_1(J)_{\langle j\rangle}=\HS_1(J_{\langle j-1\rangle})
\end{equation}
for a componentwise polymatroidal ideal $J$. This equality already fails for the principal ideal $J=(x)$ in $K[x,y]$: at $j=3$, the left side is zero, whereas
\[
 \HS_1(J_{\langle2\rangle})=\HS_1((x^2,xy))=(x^2y).
\]
The argument for \eqref{invalid} passes from least common multiples of distinct generators of a homogeneous component to least common multiples of distinct minimal generators of $J$. The former two monomials may instead be multiples of the same minimal generator of $J$. Theorem~\ref{main} shows that the final first-shift assertion itself fails, in addition to this intermediate equality. These observations concern the componentwise assertion only and do not address the equigenerated or asymptotic results of \cite{FL}.

\begin{thebibliography}{9}
\raggedright
\bibitem{CMS} Y.~Cid-Ruiz, J.~P.~Matherne, and A.~Shapiro,
\emph{Syzygies of polymatroidal ideals}, arXiv:2507.13153v1, July 17, 2025.
\url{https://arxiv.org/abs/2507.13153}.
\bibitem{F} A.~Ficarra,
\emph{Shellability of componentwise discrete polymatroids},
Electron. J. Combin. \textbf{32}(1) (2025), P1.41, published March 14, 2025.
\url{https://www.combinatorics.org/ojs/index.php/eljc/article/view/v32i1p41}.
\bibitem{FL} A.~Ficarra and D.~Lu,
\emph{Polymatroidal ideals and their asymptotic syzygies},
arXiv:2509.11977v1, September 15, 2025.
\url{https://arxiv.org/abs/2509.11977v1}.
\end{thebibliography}
\end{document}
